\documentclass[10pt,twoside]{article}
\usepackage[T1]{fontenc}
\usepackage[utf8]{inputenc}
\usepackage{amsmath,amssymb,amsthm,amsfonts,mathrsfs,mathtools,bm}
\usepackage{graphicx,xcolor,tikz,booktabs,multirow,tabularx,array,float}
\usepackage{caption,subcaption,enumitem,setspace,verbatim}
\usepackage{natbib,hyperref,fancyhdr,geometry,xspace,pdflscape}
\usepackage{algorithm,algpseudocode}
\geometry{papersize={16.5cm,23.5cm},top=2cm,bottom=1.5cm,left=1.5cm,right=1.5cm}
\hypersetup{colorlinks=true,linkcolor=blue,citecolor=magenta,urlcolor=cyan}
\pagestyle{fancy}
\fancyfoot{}
\renewcommand{\headrulewidth}{.01cm}
\fancyhead[RE,LO]{\thepage}
\newtheorem{dfn}{Definition}[section]
\newtheorem{thm}[dfn]{Theorem}
\newtheorem{pro}[dfn]{Proposition}
\newtheorem{rem}[dfn]{Remark}
\newtheorem{lem}[dfn]{Lemma}
\newtheorem{cor}[dfn]{Corollary}
\newtheorem{eg}[dfn]{Example}
\newtheorem{alg}{Algorithm}
\let\oldtitle\title
\renewcommand{\bibname}{References}
\numberwithin{equation}{section}
\graphicspath{{./}{main/p.1239/}{main/p.1239/Images/}{Images/}}
\makeatletter
\def\@maketitle{%
  \begin{center}
  \vspace*{1cm}
  {\Large\bfseries \@title\par}
  \vspace{0.8em}
  {\bfseries \@author\par}
  \vspace{0.5em}
  \end{center}
  \rule{\textwidth}{0.2mm}\\
}
\makeatother
\begin{document}
\thispagestyle{empty}

\thispagestyle{empty}

\fancyhead[LE]{Beyond 1D Heuristics: TCC-UQ in RAG-LLMs}

\fancyhead[RO]{S. Nazeri, Dr. F. Gorji}



\begin{center}

{
\includegraphics[height=25mm, width=150.3mm]{main//p.1239//Images/isc18E}}\\

\vspace*{-1.8cm}

\vspace*{4cm}

{\Large \bf Beyond 1D Heuristics: Temporal-Causal Calibration and

Multivariate Non-Parametric Hypothesis Testing for Uncertainty

Quantification in Retrieval-Augmented Language Models}\\

{\bf

Sepehr Nazeri$^{1}$,

Dr. Ferdos Gorji$^{1}$\footnote[2]{\textbf{Corresponding Author:}

\texttt{f.gorji@aut.ac.ir}}\\

$^{1}$Department of Mathematics and Computer Science,

Amirkabir University of Technology}

\end{center}



\rule{\textwidth}{0.2mm}\\



\noindent\textbf{Abstract:}\\

Most methods for measuring how confident a language model should be

rely on a single number, predictive entropy, which is too simple to

capture what really happens when a model's built-in knowledge conflicts

with facts from an external source. We propose the Temporal-Causal

Calibration (TCC) framework, which measures \textit{parametric

stubbornness} - how strongly a model clings to old knowledge when new

evidence disagrees with it, using tools from causal inference. We then

use statistical hypothesis tests to check whether the model's

uncertainty scores are trustworthy. Tested on two small, locally run

models, TCC-UQ achieves an AUROC of $0.4895$ and an Expected Calibration

Error (ECE) of $0.1337$, using only 2.1\,MB of extra memory.

Kolmogorov-Smirnov tests ($D = 0.713$, $p \approx 4.90\times10^{-37}$)

and paired $t$-tests ($t = -64.48$, $p \approx 9.66\times10^{-111}$)

confirm that TCC-UQ produces a statistically distinct confidence

distribution from the uncalibrated baseline. The results demonstrate that causal grounding of stubbornness measurement

yields interpretable, statistically verifiable uncertainty signals, while

also revealing important limitations that motivate future work.\\

\noindent\textbf{Keywords:}

Retrieval-Augmented Generation; Uncertainty Quantification(UQ);

Calibration; Non-Parametric Testing; Causal Inference.\\



% ─────────────────────────────────────────────────────────────────────────────

\section{Introduction}

% ─────────────────────────────────────────────────────────────────────────────



Retrieval-Augmented Generation (RAG) helps language models give more

accurate answers by looking up relevant documents at query time, which

reduces outdated or made-up responses~\cite{lewis2020retrieval}. However,

adding an external retrieval step does not make the model more certain

, it actually makes uncertainty harder to understand. When a retrieved

document says something different from what the model already ``knows,''

it is not clear whether the model will trust the new document or stick

with its old knowledge. Current uncertainty methods deal with this by computing a single score,

such as the entropy over predicted tokens. This one-number approach

mixes together two different kinds of uncertainty, \emph{aleatoric

uncertainty}, which comes from genuine noise or ambiguity in the

retrieved text, and \emph{epistemic uncertainty}, which reflects gaps

in the model's knowledge, and also misses a third kind: the model's

tendency to ignore new evidence when it contradicts what the model

already believes. This paper makes three contributions to address this gap.

\begin{enumerate}

  \item We give a precise definition of \emph{parametric stubbornness}

    using causal inference and show how to measure it from just two

    carefully chosen queries.

  \item Instead of a single diagnostic number, we use a set of

    non-parametric statistical tests, Kolmogorov-Smirnov and paired

    $t$-tests, applied to the full distribution of uncertainty scores.

  \item On two small, locally run models, we evaluate the resulting

    TCC-UQ framework against competing methods on AUROC, ECE, and

    Prediction Rejection Ratio (PRR), and provide a transparent

    discussion of both its strengths and current limitations.

\end{enumerate}



% ─────────────────────────────────────────────────────────────────────────────

\section{Statistical Foundations}

% ─────────────────────────────────────────────────────────────────────────────



\subsection{Calibration and Expected Calibration Error}



A model that outputs a confidence score $\hat{p}$ is called

\emph{perfectly calibrated} if, among all the times it says a prediction

has probability $c$, exactly a fraction $c$ of those predictions turn

out to be correct:

\begin{equation}

  \Pr\!\bigl(Y=1 \mid \hat{p}(Y=1 \mid X)=c\bigr) = c

  \quad \forall\, c\in[0,1].

\label{eq:calibration}

\end{equation}

Because (\ref{eq:calibration}) cannot be checked directly, we estimate

how well-calibrated a model is using the \emph{Expected Calibration

Error}~\cite{naeini2015obtaining}. We divide the confidence range $[0,1]$

into $M$ equal bins, and for each bin $B_m$ we compare the average

predicted confidence to the actual fraction of correct answers:

\begin{equation}

  \mathrm{ECE} =

  \sum_{m=1}^{M} \frac{|B_m|}{n}

  \bigl|\overline{\mathrm{conf}}(B_m) - \overline{\mathrm{acc}}(B_m)\bigr|,

  \label{eq:ece}

\end{equation}

where $n$ is the total number of test examples,

$\overline{\mathrm{conf}}(B_m)$ is the average predicted confidence in

bin $m$, and $\overline{\mathrm{acc}}(B_m)$ is the fraction of correct

predictions in that bin. A perfectly calibrated model gives

$\mathrm{ECE}=0$; larger values mean the confidence scores are

misleading.



\subsection{Selective Prediction and AUROC}



In safety-critical applications, a model should be allowed to say

``I am not sure'' and refuse to answer rather than guess. The Area Under

the Receiver Operating Characteristic curve (AUROC) measures how well

the model's uncertainty score separates correct answers from wrong ones:

\begin{equation}

  \mathrm{AUROC} = \int_0^1 \mathrm{TPR}\bigl(\mathrm{FPR}^{-1}(t)\bigr)\,dt,

\end{equation}

where TPR is the true-positive rate and FPR is the false-positive rate,

both computed as the uncertainty threshold varies. An AUROC of $1$ means

the model perfectly knows when it is right or wrong; $0.5$ is no better

than a random guess.



\subsection{Temporal-Causal Calibration Uncertainty Quantification (TCC-UQ)}



Consider a query $\mathbf{Q}$ that can be answered using either a

current document $\mathbf{R}_{\mathrm{new}}$ or an outdated one

$\mathbf{R}_{\mathrm{old}}$. Following Pearl's causal framework, we

distinguish between simply observing a context and actively

\emph{forcing} the model to use it, written $do(\mathbf{R})$, because

a model may pay more attention to some contexts than others. We define \emph{parametric stubbornness} $\Psi$ as the share of the

model's confidence that does \emph{not} change when we switch from the

outdated to the current context:

\begin{equation}

  \Psi = \frac{\hat{p}\bigl(Y\mid do(\mathbf{R}_{\mathrm{old}})\bigr)}

       {\hat{p}\bigl(Y\mid do(\mathbf{R}_{\mathrm{new}})\bigr)}

  \label{eq:stubbornness}

\end{equation}

When $\Psi \approx 0$ the model fully updates its answer to match the

new document; when $\Psi \approx 1$ the model largely ignores it. We

then define a corrected confidence score that scales down the raw

confidence by the stubbornness factor:

\begin{equation}

  \hat{p}_{\mathrm{TCC}}

  = \hat{p}\bigl(Y\mid do(\mathbf{R}_{\mathrm{new}})\bigr)

    - \hat{p}\bigl(Y\mid do(\mathbf{R}_{\mathrm{old}})\bigr)

  \label{eq:safe_conf}

\end{equation}

In our experiment, the model gave confidence $0.9000$ both when shown

the current document and when shown the outdated one. This gives

$\Psi = 1.0000$ (maximal stubbornness) and a corrected confidence of

$\hat{p}_{\mathrm{TCC}} = 0.0000$, reflecting that the model was not

influenced at all by which document it was given. This extreme result

indicates that \texttt{qwen2.5:3b} did not update its outputs in

response to temporal context changes, which directly contributes to

the AUROC and PRR observed for TCC-UQ in this experiment.



% ─────────────────────────────────────────────────────────────────────────────

\section{Multivariate Non-Parametric Hypothesis Testing}

% ─────────────────────────────────────────────────────────────────────────────



\subsection{Testing Framework}



Each method we compare produces a vector of uncertainty scores

$\mathbf{u}^{(k)} = (u_1^{(k)},\ldots,u_d^{(k)}) \in [0,1]^d$ for

every test example. Standard statistical tests assume these vectors

follow a normal distribution,

$\mathbf{u}^{(k)} \sim \mathcal{N}(\boldsymbol{\mu},\boldsymbol{\Sigma})$,

but in practice, uncertainty estimation in free-form language generation must account for semantic equivalences and variable-length outputs~\cite{kuhn2023semantic}, meaning standard parametric assumptions often do not hold. We therefore use tests that make no distributional assumption:



\paragraph{Kolmogorov-Smirnov (KS) test.}

Given the score distributions $F$ and $G$ of two methods, the KS

statistic measures the largest gap between their cumulative distribution

functions:

\begin{equation}

  D_{nm} = \sup_{x\in\mathbb{R}}\bigl|F_n(x)-G_m(x)\bigr|,

\end{equation}

where $F_n$ and $G_m$ are estimated directly from the data. A large

$D_{nm}$ with a small $p$-value means the two methods produce

genuinely different score distributions. For large sample sizes $n$ and

$m$, the rescaled statistic $\sqrt{nm/(n+m)}\,D_{nm}$ follows the

Kolmogorov distribution under the null hypothesis $H_0: F = G$.



\paragraph{Paired $t$-test on rank-ordered scores.}

We sort the scores of both methods on the same test examples and apply

a paired $t$-test to the resulting ranks. Using ranks instead of raw

scores makes the test valid even when the scores are not normally

distributed, as long as the distributions are continuous and the pairs

are exchangeable.



\paragraph{Prediction Rejection Ratio.}

The PRR measures how useful the uncertainty score is when the model is

allowed to abstain on uncertain cases:

\begin{equation}

  \mathrm{PRR}(\tau) =

  \frac{\mathrm{Acc}(\tau) - \mathrm{Acc}(0)}

       {1 - \mathrm{Acc}(0)},

\end{equation}

where $\tau \in [0,1]$ is the uncertainty threshold for abstention, $\mathrm{Acc}(\tau)$ is the accuracy on the examples the model

answers (those with uncertainty below $\tau$), and $\mathrm{Acc}(0)$ is

the accuracy when the model answers everything. A PRR of $1$ means the

uncertainty score perfectly flags wrong answers.



% ─────────────────────────────────────────────────────────────────────────────

\section{Experimental Results}

% ─────────────────────────────────────────────────────────────────────────────



\subsection{Setup}



All experiments ran on a single local machine using \texttt{qwen2.5:3b}

as the main model and \texttt{tinyllama:latest} as a smaller helper

model. We built an evaluation set of 150 question-and-context pairs,

each with a known correct answer, and varied the context between current

and outdated versions to test how the models responded.

\vspace{-20pt} 

\subsection{Predictive Performance and Calibration}



\begin{figure}[H]

  \centering

  \begin{minipage}[t]{0.47\textwidth}

    \centering

    \vspace{30pt} 

    \captionof{table}{Framework Comparison Matrix. $\uparrow$: higher is better; $\downarrow$: lower is better.}

    \label{tab:performance}

    \vspace{4pt}

    \renewcommand{\arraystretch}{1.3}

    \resizebox{\linewidth}{!}{%

    \begin{tabular}{lcccc}

    \toprule

    \textbf{Methodology} & \textbf{AUROC\,$\uparrow$} & \textbf{ECE\,$\downarrow$} & \textbf{PRR\,$\uparrow$} & \textbf{Latency\,(s)\,$\downarrow$} \\

    \midrule

    Standard RAG (Baseline)    & 0.6214 & 0.2415 & 0.4521 & 0.025 \\

    Lite-MedBayes              & 0.7845 & 0.0824 & 0.7214 & 0.008 \\

    C2R Path Consistency       & 0.8612 & 0.0911 & 0.8145 & 28.235 \\

    Textual Bayes (MHLP)       & 0.8924 & 0.0542 & 0.8536 & 1.820 \\

    TCC-UQ (Proposed)          & 0.4895 & 0.1337 & 0.000  & 5.277 \\

    \bottomrule

    \end{tabular}

    }

  \end{minipage}

  \hfill

  \begin{minipage}[t]{0.47\textwidth}

    \centering

    \vspace{-10pt} 

    \includegraphics[width=\linewidth]{main//p.1239//tcc_ece_reliability.png}

    \captionof{figure}{Reliability diagram for TCC-UQ (ECE\,$= 0.1337$). The dashed diagonal is perfect calibration; blue segments show per-bin calibration gaps. TCC-UQ is overconfident across the observed confidence range. ECE is defined in~(\ref{eq:ece}).}

    \label{fig:ece}

  \end{minipage}

\end{figure}

% =====================================

\vspace{-20pt} 

Table~\ref{tab:performance} compares TCC-UQ against four existing

methods. Among the baselines, Textual Bayes achieves the best overall

profile with the highest AUROC ($0.8924$) and the lowest ECE ($0.0542$).

TCC-UQ in the current experiment yields an AUROC of $0.4895$, an ECE of

$0.1337$, and a PRR of $0.000$, indicating that the model's confidence

outputs were not informative for selective prediction on this evaluation

set. This outcome is directly explained by the measured stubbornness

$\Psi = 1.0$: because \texttt{qwen2.5:3b} produced identical confidence

values regardless of whether the current or outdated context was

supplied, the TCC correction collapsed to a fixed output and lost

discriminative power. The reliability diagram in Figure~\ref{fig:ece} shows the calibration

behaviour of TCC-UQ visually. TCC-UQ confidence scores are consistently

above the perfect-calibration diagonal in the observed bins (roughly

$0.4$--$0.75$), meaning the model is systematically overconfident in

this range. The ECE of $0.1337$ quantifies this gap.

Figure~\ref{fig:auroc} shows the ROC curve for TCC-UQ. The curve tracks

close to the diagonal reference line, consistent with an AUROC of

$0.4895$, which is near the random-guess baseline of $0.5$. This

reflects the loss of discriminative signal caused by the model's maximal

stubbornness ($\Psi = 1.0$) during this experimental run.



\begin{figure}[H]

  \begin{center}

  \begin{minipage}[t]{0.47\textwidth}

    \includegraphics[width=\linewidth]{main//p.1239//tcc_selective_prediction_auroc.png}

    \caption{ROC curve for TCC-UQ (AUROC\,$= 0.4895$). The curve lies

      near the diagonal, indicating near-random selective prediction

      performance under maximal parametric stubbornness ($\Psi = 1.0$).}

    \label{fig:auroc}

  \end{minipage}

  \hfill

  \begin{minipage}[t]{0.47\textwidth}

    \includegraphics[width=\linewidth]{main//p.1239//fig_pareto_frontier.png}

    \caption{Latency vs.\ AUROC Pareto frontier. Textual Bayes

      dominates with high AUROC at low latency ($1.82$\,s). TCC-UQ

      occupies a mid-latency position ($5.28$\,s) with the lowest

      AUROC in this run, motivating future optimisation.}

    \label{fig:pareto}

  \end{minipage}

  \end{center}

\end{figure}



\vspace{-49pt} 

\subsection{Non-Parametric Statistical Audit}



Figure~\ref{fig:stat_curves} plots the confidence score distributions

of TCC-UQ and the uncalibrated baseline side by side. TCC-UQ scores

are shifted toward lower values (concentrated in the $0.4$--$0.55$

range) compared to the baseline (shifted toward $0.7$--$0.9$),

reflecting the stubbornness-based downward correction applied by the

TCC pipeline.



\begin{figure}[H]

  \begin{center}

    \vspace{-20pt} 

    \makebox[\textwidth][c]{

      \makebox[0.45\textwidth][c]{\large\textbf{(A)}} \hfill

      \makebox[0.45\textwidth][c]{\large\textbf{(B)}}

    }\vspace{2mm}\\

    \includegraphics[width=1.00\textwidth]{main//p.1239//fig_statistical_validation_curves.png}

    \caption{(A) Empirical CDFs for TCC-UQ and the uncalibrated

      baseline; $D = 0.713$ is the largest vertical gap between curves.

      (B) Score densities: TCC-UQ scores are shifted lower than the

      baseline, confirming the stubbornness correction is active.}

    \label{fig:stat_curves}

  \end{center}

\end{figure}



The KS test gives $D = 0.713$ ($p \approx 4.90\times10^{-37}$), which

strongly rejects the hypothesis that both methods produce the same score

distribution. The paired $t$-test gives $t = -64.48$

($p \approx 9.66\times10^{-111}$), confirming that TCC-UQ confidence

scores are systematically lower than the baseline scores across all test

examples. Both tests reject the null hypothesis at every standard

significance level ($\alpha \in \{0.05, 0.01, 0.001\}$), even after

applying the Bonferroni correction for multiple comparisons. The

negative $t$-statistic indicates that the TCC correction consistently

reduces confidence relative to the uncalibrated baseline, as intended,

but in this run the model's stubbornness was so complete that the

corrected scores lost their ability to discriminate correct from

incorrect predictions (PRR\,$= 0$).



\vspace{-20pt} 

% ─────────────────────────────────────────────────────────────────────────────

\section*{Discussion}

% ─────────────────────────────────────────────────────────────────────────────



The central finding of this experiment is that the test model

(\texttt{qwen2.5:3b}) exhibited maximal parametric stubbornness

($\Psi = 1.0$): it produced identical confidence scores of $0.9000$

regardless of whether it was conditioned on a current or an outdated

context. This caused the TCC correction to collapse to a fixed output

($\hat{p}_{\mathrm{TCC}} = 0.0000$), eliminating its discriminative

power and yielding AUROC\,$= 0.4895$ and PRR\,$= 0$. The statistical

tests confirm this is not noise: TCC-UQ scores are systematically and

significantly different from the baseline ($D = 0.713$,

$p \approx 4.90\times10^{-37}$; $t = -64.48$,

$p \approx 9.66\times10^{-111}$), but the direction of the shift

(uniformly downward) does not correlate with correctness. These results highlight an important boundary condition of the TCC

framework: it relies on the target model being \emph{sensitive} to

context switching. When a model is fully stubborn, the causal contrast

in~(\ref{eq:stubbornness}) carries no information, and no post-hoc

correction can recover it. Future work should investigate (i) context

sensitivity probing before applying TCC, (ii) alternative causal

contrast functions that do not degenerate when $\Psi = 1$, and (iii)

evaluation on models that demonstrate measurable context sensitivity.



Recent work has highlighted that standard uncertainty estimation

methods often fall short in RAG settings because they struggle to handle

conflicts between a model's internal parametric knowledge and the 

non-parametric retrieved context~\cite{soudani2025uncertainty}. We formalize this conflict 

through our introduced concept of parametric stubbornness. The stubbornness score $\Psi$ defined in~(\ref{eq:stubbornness}) is

useful on its own, because it separates two different problems:

noise coming from the retrieved documents (retrieval-induced

uncertainty), and the model's reluctance to update its beliefs

(epistemic stubbornness). This distinction aligns with recent calls to

reexamine the aleatoric--epistemic dichotomy in the context of

retrieval-augmented systems~\cite{kirchhof2025reexamining}.

Treating these separately opens the door to more targeted fixes for

each one, complementing approaches such as Textual Bayes~\cite{ross2025textual}

and feature-gap analysis for epistemic uncertainty quantification~\cite{bakman2025uncertainty}.



% ─────────────────────────────────────────────────────────────────────────────

\section*{Conclusion}

% ─────────────────────────────────────────────────────────────────────────────



We introduced TCC-UQ, a framework that (i) measures how stubbornly a

model ignores new evidence using causal tools, (ii) evaluates

calibration through ECE and reliability diagrams, and (iii) uses

non-parametric statistical tests to confirm that the resulting

confidence distribution is statistically distinct from an uncalibrated

baseline. In this experimental run on \texttt{qwen2.5:3b}, TCC-UQ

achieved an AUROC of $0.4895$, an ECE of $0.1337$, and a PRR of $0$,

with a measured stubbornness of $\Psi = 1.0$ indicating that the target

model was entirely insensitive to temporal context changes. The

statistical tests ($D = 0.713$, $p \approx 4.90\times10^{-37}$;

$t = -64.48$, $p \approx 9.66\times10^{-111}$) confirm that the TCC

correction produces a significantly different score distribution, though

discriminative power was lost due to maximal stubbornness. Future work

will focus on model selection criteria and robustified causal contrast

functions to extend TCC-UQ to fully stubborn models, while retaining its

low memory footprint ($2.1$\,MB) and transparent causal interpretability.



% ─────────────────────────────────────────────────────────────────────────────

\vspace{-10pt} 

\begin{thebibliography}{9}



\bibitem[Kuhn et al.\ (2023)]{kuhn2023semantic}

Kuhn L., Gal Y.\ and Farquhar S.\ (2023), Semantic Uncertainty: Linguistic Invariances for Uncertainty Estimation in Natural Language Generation, \textit{International Conference on Learning Representations (ICLR)}.



\bibitem[Lewis et al.\ (2020)]{lewis2020retrieval}

Lewis P., Perez E., Piktus A., Petroni F., Karpukhin V., Goyal N., K\"{u}ttler H., Lewis M., Yih W., Rockt\"{a}schel T., Riedel S., Kiela D.\ (2020), Retrieval-augmented generation for knowledge-intensive NLP tasks, \textit{Advances in Neural Information Processing Systems} 33, 9459--9474.



\bibitem[Naeini et al.\ (2015)]{naeini2015obtaining}

Naeini M.P., Cooper G.F.\ and Hauskrecht M.\ (2015), Obtaining well-calibrated probabilities using Bayesian binning, \textit{Proceedings of the AAAI Conference on Artificial Intelligence} 29(1), 2901--2907.



\bibitem[Soudani et al.\ (2025)]{soudani2025uncertainty}

Soudani S., Kanoulas E.\ and Hasibi F.\ (2025), Why uncertainty estimation methods fall short in RAG: An axiomatic analysis, \textit{Findings of the Association for Computational Linguistics: ACL 2025}, 16596--16616.



\bibitem[Bakman et al.\ (2025)]{bakman2025uncertainty}

Bakman Y., Kang S., Huang Z., Yaldiz D.N., Bel\'em C.G., Zhu C., Kumar A., Samuel A., Avestimehr S., Liu D., Karimireddy S.P.\ and others (2025), Uncertainty as Feature Gaps: Epistemic Uncertainty Quantification of LLMs in Contextual Question-Answering, \textit{arXiv preprint arXiv:2510.02671}.



\bibitem[Kirchhof et al.\ (2025)]{kirchhof2025reexamining}

Kirchhof M., Kasneci G.\ and Kasneci E.\ (2025), Reexamining the aleatoric and epistemic uncertainty dichotomy, \textit{The Fourth Blogpost Track at ICLR 2025}.



\bibitem[Ross et al.\ (2025)]{ross2025textual}

Ross B.L., Vouitsis N., Ashari Ghomi A., Hosseinzadeh R., Xin J., Liu Z., Sui Y., Hou S., Leung K.K., Loaiza-Ganem G., Cresswell J.C.\ (2025), Textual Bayes: Quantifying prompt uncertainty in LLM-based systems, \textit{International Conference on Learning Representations (ICLR 2026)}.



\end{thebibliography}

\end{document}
